
\documentclass[11pt]{article}
\usepackage{series}
\usepackage{hyperref}
\usepackage{booktabs}
\usepackage{array}

\title{Path Integral Constraints as Finite Admissibility Filters\\
\large Delta Functionals, Soft Constraints, and Projection in Modal Triplet Theory}
\author{Peter Nero}
\date{April 2026}

\begin{document}
\maketitle

\begin{abstract}
Delta functionals occur throughout path-integral physics: as hard constraints, gauge-fixing conditions,
phase-space reductions, Lagrange-multiplier insertions, and exact selection of histories.  In standard
formalism they are usually treated as distributional devices.  This paper continues the delta/projection
program by showing that a delta functional may be understood as the zero-width limit of a finite
admissibility filter over histories.  The rigorous core is finite-dimensional: for a smooth constraint map
\(C:M\to\mathbb R^m\) with \(0\) a regular value, normalized Gaussian filters
\[
\mathcal K_\varepsilon(C(x))=(2\pi\varepsilon^2)^{-m/2}\exp(-|C(x)|^2/2\varepsilon^2)
\]
converge, as measures, to the surface delta \(\delta(C(x))\,dx\), including the standard Jacobian factor
from the coarea formula.  The infinite-dimensional path-integral statement is then treated as a formal
extension requiring regularization.  In MTT language, the hard delta \(\delta[C]\) is the singular
limit of a finite admissibility weight \(\exp[-J_{\mathrm{adm}}/\epsilon^2]\).  This also clarifies the relation between multiplier representations, large-penalty localization, and finite admissibility tubes, and unifies the
interpretation of classical constraint shells, gauge-fixing deltas, and path-integral restrictions:
exact constraint enforcement is the zero-width idealization of finite admissibility selection.
\end{abstract}

\tableofcontents

\section{Purpose and claim discipline}

The previous papers in this sequence reinterpreted several recurring delta functions:
identity kernels, point sources, gauge slices, measurement projectors, contact vertices,
conservation deltas, spectral peaks, white noise, and microcanonical shells.
The present paper treats the path-integral version:
\[
\delta[C[\phi]] .
\]

The guiding claim is:
\[
\boxed{\delta[C]=\text{zero-width limit of a finite admissibility filter over histories}.}
\]

This paper has two layers.

\begin{center}
\begin{tabular}{p{0.28\linewidth}p{0.62\linewidth}}
\toprule
Layer & Status \\
\midrule
Finite-dimensional constraint theorem & rigorous, proved using the coarea formula \\
Functional-integral extension & formal; requires regulator, measure choice, and anomaly control \\
MTT interpretation & structural; finite admissibility weights replace hard delta constraints \\
\bottomrule
\end{tabular}
\end{center}

The paper does not claim to define all path integrals nonperturbatively.  It isolates one
mathematical pattern: hard functional deltas are singular limits of finite-width constraint
filters.

\section{From hard constraints to finite filters}

A standard constrained path integral has the schematic form
\[
Z=\int \mathcal D\phi\; \delta[C[\phi]]\,e^{iS[\phi]/\hbar}.
\]
The delta functional enforces exact membership in the constraint surface \(C[\phi]=0\).

The finite-filter replacement is
\[
\delta[C[\phi]]
\quad\leadsto\quad
\mathcal K_\varepsilon[C[\phi]]
=
\exp\!\left(-\frac{1}{2\varepsilon^2}\|C[\phi]\|^2\right),
\]
possibly with normalization and with the norm chosen by an admissibility metric.
The finite-filter path integral is then
\[
Z_\varepsilon
=
\int \mathcal D\phi\;
\mathcal K_\varepsilon[C[\phi]]\,
e^{iS[\phi]/\hbar}.
\]

In Euclidean signature one may instead write
\[
Z_\varepsilon^{E}
=
\int \mathcal D\phi\;
\exp\!\left(
-S_E[\phi]
-\frac{1}{2\varepsilon^2}\|C[\phi]\|^2
\right).
\]
This is a soft-constraint or penalty formulation.  The hard constraint is recovered only in
the \(\varepsilon\downarrow0\) limit.

\section{Finite-dimensional theorem}

Let \(M\) be an \(n\)-dimensional smooth Riemannian manifold with volume measure \(d\mu\).
Let
\[
C:M\to\mathbb R^m,\qquad m\le n,
\]
be a smooth constraint map, and assume \(0\) is a regular value.  Then
\[
\Sigma:=C^{-1}(0)
\]
is a smooth submanifold of codimension \(m\).

Define the constraint Jacobian
\[
J_C(x):=\sqrt{\det\bigl(DC_x\,DC_x^{\ast}\bigr)},
\]
where \(DC_x:T_xM\to\mathbb R^m\) and the adjoint is taken with respect to the Riemannian
metric on \(M\) and the Euclidean metric on \(\mathbb R^m\).

The surface delta is defined by
\[
\int_M f(x)\,\delta(C(x))\,d\mu(x)
=
\int_{\Sigma}\frac{f(x)}{J_C(x)}\,d\sigma(x),
\]
for test functions \(f\), where \(d\sigma\) is the induced hypersurface measure.

\subsection{Gaussian constraint filters}

Let
\[
k_\varepsilon(u):=(2\pi\varepsilon^2)^{-m/2}\exp(-|u|^2/2\varepsilon^2).
\]

\begin{theorem}[Gaussian constraint filters converge to surface deltas]
Let \(f\in C_c^\infty(M)\).  Under the regular-value assumption above,
\[
\lim_{\varepsilon\downarrow0}
\int_M f(x) k_\varepsilon(C(x))\,d\mu(x)
=
\int_{\Sigma}\frac{f(x)}{J_C(x)}\,d\sigma(x).
\]
Equivalently,
\[
k_\varepsilon(C(x))\,d\mu(x)\;\longrightarrow\;\delta(C(x))\,d\mu(x)
\]
weakly as measures.
\end{theorem}

\begin{proof}
By the coarea formula, for any integrable \(h\),
\[
\int_M h(x)J_C(x)\,d\mu(x)
=
\int_{\mathbb R^m}
\left(\int_{C^{-1}(u)} h(x)\,d\sigma_u(x)\right)\,du .
\]
Apply this with
\[
h(x)=\frac{f(x)k_\varepsilon(C(x))}{J_C(x)} ,
\]
which is well-defined on a neighborhood of \(\Sigma\) because \(0\) is a regular value and
\(f\) has compact support.  Then
\[
\int_M f(x)k_\varepsilon(C(x))\,d\mu(x)
=
\int_{\mathbb R^m} k_\varepsilon(u)
\left[
\int_{C^{-1}(u)}
\frac{f(x)}{J_C(x)}\,d\sigma_u(x)
\right]du .
\]
For \(u\) near \(0\), the bracketed expression is smooth by the regular-level-set theorem
and compact support of \(f\).  Denote it by \(F(u)\).  Since \(k_\varepsilon\) is an approximate
identity on \(\mathbb R^m\),
\[
\int_{\mathbb R^m} k_\varepsilon(u)F(u)\,du\to F(0).
\]
But
\[
F(0)=\int_{\Sigma}\frac{f(x)}{J_C(x)}\,d\sigma(x).
\]
This proves the claim.
\end{proof}

\section{General admissibility kernels}

Gaussian filters are convenient but not essential.  Let
\[
k_\varepsilon(u)=\varepsilon^{-m}k(u/\varepsilon),
\]
where \(k\in C_c^\infty(\mathbb R^m)\), \(k\ge0\), and \(\int k(u)\,du=1\).
The same proof gives:
\[
k_\varepsilon(C(x))\,d\mu(x)\to \delta(C(x))\,d\mu(x).
\]

Thus the hard delta does not depend on the detailed shape of the finite filter.  The detailed
shape controls finite-\(\varepsilon\) corrections.  In MTT language, those corrections encode
finite admissibility width.

\section{Path integrals: formal functional extension}

The formal field-theoretic analogue is obtained by replacing \(x\in M\) with a field
configuration \(\phi\in\mathcal F\) and \(C(x)\) with a functional constraint
\[
C[\phi]\in \mathcal Y .
\]
A hard constrained functional integral has the schematic form
\[
Z=
\int_{\mathcal F}\mathcal D\phi\;
\delta(C[\phi])\,e^{iS[\phi]/\hbar}.
\]

The finite-filter version is
\[
Z_\varepsilon=
\int_{\mathcal F}\mathcal D\phi\;
\mathcal K_\varepsilon(C[\phi])\,e^{iS[\phi]/\hbar},
\]
where
\[
\mathcal K_\varepsilon(C)
=
\mathcal N_\varepsilon
\exp\!\left(-\frac{1}{2\varepsilon^2}\|C\|_{\mathcal Y}^{2}\right).
\]

In Euclidean form:
\[
Z_\varepsilon^E
=
\int_{\mathcal F}\mathcal D\phi\;
\exp\!\left(
-S_E[\phi]
-\frac{1}{2\varepsilon^2}\|C[\phi]\|_{\mathcal Y}^{2}
\right).
\]

\subsection{Scope warning}

In infinite dimensions, the symbols
\[
\mathcal D\phi,\qquad \delta(C[\phi]),\qquad \det(DC\,DC^\ast)
\]
are formal unless a regulator, Gaussian measure, lattice approximation, pAQFT construction,
or other definition is supplied.  The rigorous theorem of this paper is finite-dimensional.
The functional-integral formulas are structural extensions of the same mechanism.

\section{Relation to Lagrange multipliers}

A delta constraint may also be represented by a Fourier integral:
\[
\delta(C)=\int \frac{d\lambda}{2\pi}\,e^{i\lambda C}.
\]
For multiple constraints:
\[
\delta^{(m)}(C)=\int \frac{d^m\lambda}{(2\pi)^m}\,e^{i\lambda\cdot C}.
\]

This expresses hard constraint enforcement by an auxiliary multiplier field.  The finite-filter
version corresponds to adding a damping envelope in multiplier space:
\[
k_\varepsilon(C)
=
\int \frac{d^m\lambda}{(2\pi)^m}
e^{i\lambda\cdot C}\,
e^{-\varepsilon^2|\lambda|^2/2}.
\]

Thus finite admissibility width in constraint space is equivalent to suppressing arbitrarily
large Lagrange-multiplier modes.


\section{Large-penalty localization}

The same limit may be expressed dynamically as a large-penalty limit.  In Euclidean
signature, replace the constrained integral
\[
\int \delta(C(x))\,e^{-S(x)}\,d\mu(x)
\]
by
\[
\int
\exp\!\left[-S(x)-\frac{|C(x)|^2}{2\varepsilon^2}\right]\,d\mu(x).
\]
As \(\varepsilon\downarrow0\), the penalty term suppresses all configurations outside an
\(O(\varepsilon)\)-tube around the constraint surface \(C^{-1}(0)\).  In local coordinates
\((u,v)\), where \(u=C(x)\in\mathbb R^m\) and \(v\) parametrizes the constraint surface,
the integral has the schematic form
\[
\int e^{-|u|^2/2\varepsilon^2}\,F(u,v)\,du\,dv .
\]
After multiplication by the normalizing factor \((2\pi\varepsilon^2)^{-m/2}\), the Gaussian
integral localizes to \(u=0\), giving
\[
\int_{C^{-1}(0)}
\frac{e^{-S(x)}}{J_C(x)}\,d\sigma(x),
\]
with \(J_C=(\det DC\,DC^\ast)^{1/2}\).  Thus the finite admissibility filter is not merely
a visual analogy for a delta; it has the same localizing asymptotics as the hard constrained
measure.

In Lorentzian signature the corresponding statement is more delicate because the path
integral is oscillatory.  The penalty term may be introduced after Wick rotation, by
\(i0\)-type damping, or inside a regulator-defined construction.  The structural point remains
the same: the hard functional delta is the zero-width limit of a filter that increasingly
localizes histories near \(C=0\).

\section{Admissible filters versus arbitrary penalties}

A finite-width constraint filter changes the theory at finite \(\varepsilon\).  Therefore the
replacement
\[
\delta[C]\leadsto \mathcal K_\varepsilon[C]
\]
is not automatically legitimate as physics.  It is legitimate only if the filter is either:

\begin{enumerate}
\item a regulator that is removed in the end;
\item an effective model for finite experimental or coarse-graining resolution; or
\item an admissibility kernel derived from the underlying MTT projection and closure-cost
structure.
\end{enumerate}

This distinction is important.  A generic penalty may violate a symmetry, distort the measure,
or select the wrong sector.  In gauge theory, for example, finite gauge-tube filters must be
compatible with the Faddeev--Popov/BV--BRST quotient structure if they are to define the
same physical quotient in the sharp limit.  In MTT terms, the filter must be an admissible
projection weight, not an arbitrary smearing inserted by hand.


\section{Relation to gauge fixing}

Gauge fixing is a special case of constraint filtering.  A gauge condition
\[
G[A]=0
\]
inserts
\[
\delta(G[A])
\]
together with the Faddeev--Popov determinant.

The finite admissibility version is
\[
\delta(G[A])
\quad\leadsto\quad
\mathcal K_\varepsilon(G[A])
=
\exp\!\left(-\frac{1}{2\varepsilon^2}\|G[A]\|^2\right).
\]

This replaces a hard gauge slice by a finite gauge tube.  The Faddeev--Popov determinant
remains the Jacobian of projection from gauge-orbit coordinates to the chosen slice, while
the filter width controls how sharply the representative section is enforced.

This paper therefore generalizes the gauge paper:
\[
\text{gauge-fixing delta}
\subset
\text{path-integral constraint delta}
\subset
\text{finite admissibility filter limit}.
\]

\section{MTT interpretation}

In Modal Triplet Theory, admissible descriptions are not arbitrary configurations.  They are
configurations that survive projection, stabilization, and coherence constraints.  Therefore
a hard functional delta
\[
\delta(C[\phi])
\]
should not be treated as primitive.  It is the zero-width limit of an admissibility filter:
\[
\mathcal K_{\epsilon}[C[\phi]]
=
\exp\!\left(-\frac{1}{\epsilon^2}J_{\mathrm{adm}}[\phi]\right),
\]
where \(J_{\mathrm{adm}}\) measures the degree of constraint violation, closure strain, gauge-slice
distance, or failure of coherent continuation.

The hard limit is
\[
\delta(C[\phi])
=
\lim_{\epsilon\downarrow0}\mathcal K_\epsilon[C[\phi]]
\]
in the same structural sense that finite-dimensional Gaussian filters converge to surface
deltas.

\subsection{Triadic reading}

The same delta may involve different MTT carrier roles.

\begin{center}
\begin{tabular}{p{0.22\linewidth}p{0.68\linewidth}}
\toprule
Carrier role & Constraint-filter interpretation \\
\midrule
Circle \(C\) & exact bookkeeping closure, return consistency, conservation shells \\
Lens \(L\) & representative selection inside redundancy classes, gauge slices \\
Nil \(N\) & survivor-basin thresholds, admissibility boundaries, collapse-like selection \\
\bottomrule
\end{tabular}
\end{center}

A path-integral delta functional is therefore not a single primitive object in MTT.  It is a
downstream singular notation for whichever admissibility condition is being enforced.

\section{Soft constraints versus hard constraints}

The finite filter
\[
\exp(-\|C\|^2/2\varepsilon^2)
\]
does not merely approximate a mathematical delta.  For finite \(\varepsilon\), it defines a different
effective theory: one in which constraint violation is penalized rather than forbidden.

This distinction matters.

\begin{center}
\begin{tabular}{p{0.28\linewidth}p{0.32\linewidth}p{0.30\linewidth}}
\toprule
Regime & Mathematical form & MTT reading \\
\midrule
Hard constraint & \(\delta(C)\) & zero-width admissibility limit \\
Soft constraint & \(e^{-\|C\|^2/2\varepsilon^2}\) & finite admissibility tolerance \\
Penalty action & \(S+\|C\|^2/2\varepsilon^2\) & closure-cost deformation \\
Unconstrained & no filter & no enforced admissibility shell \\
\bottomrule
\end{tabular}
\end{center}

The hard limit may be useful and often correct.  But MTT treats it as a limit, not as the
primitive description.

\section{Diagnostics and finite-width effects}

Replacing \(\delta[C]\) by \(\mathcal K_\varepsilon[C]\) suggests concrete finite-width diagnostics:

\begin{enumerate}
\item constraint surfaces become finite admissibility tubes;
\item gauge slices become finite gauge tubes;
\item hard classical shells become soft shells;
\item path histories with small closure violation contribute with controlled suppression;
\item infinitely sharp Lagrange-multiplier modes are damped;
\item exact constraint enforcement is recovered only as \(\varepsilon\downarrow0\).
\end{enumerate}

The physical relevance of these corrections depends on the sector.  In ordinary effective
field theory they may be only regulator artifacts.  In MTT they may encode finite coherence
capacity.

\section{Scope of proof}

\begin{center}
\begin{tabular}{p{0.25\linewidth}p{0.65\linewidth}}
\toprule
Statement & Status \\
\midrule
Gaussian filters on regular finite-dimensional constraint surfaces converge to surface deltas & proved \\
General approximate-identity filters give the same surface delta & proved by same coarea argument \\
Large-penalty filters localize to the constraint surface in the sharp limit & proved locally by the same coarea/Gaussian argument \\
Functional delta filters behave analogously in path integrals & formal, regulator-dependent \\
Gauge-fixing deltas are a special case of constraint filtering & structural and standard \\
Finite penalties preserve physics only when regulator-removable or admissibility-derived & caveat, not automatic \\
MTT admissibility costs generate all physical constraint filters & interpretive target, not proved here \\
\bottomrule
\end{tabular}
\end{center}

\section{Conclusion}

Path-integral delta functionals are usually written as hard constraints:
\[
\delta[C[\phi]].
\]
This paper shows that, at the rigorous finite-dimensional level, such deltas arise as
zero-width limits of finite constraint filters.  The same pattern extends formally to functional
integrals, where a hard delta can be replaced by a finite admissibility weight
\[
\exp[-\|C[\phi]\|^2/2\varepsilon^2].
\]

The MTT reading is direct:
\[
\boxed{
\delta[C]=\text{singular shadow of finite admissibility filtering}.
}
\]
This unifies classical constraint shells, gauge-fixing deltas, soft penalty actions, and
path-integral restrictions under one projection-first interpretation.

\end{document}
